\documentclass[11pt,a4paper]{article}
\usepackage[margin=25mm]{geometry}

\title{Window Length and Sydney Transport Mode Classification}
\author{}
\date{}

\begin{document}
\maketitle

\begin{abstract}
This study compares 12 window lengths from 20 to 240 seconds for Sydney bus,
car, and train classification. A fixed XGBoost protocol evaluates the same
149 sessions and seven held-out participants at every length. Mean pooled
macro F1 across three seeds peaks at 0.693 for 120-second windows. The result
is descriptive: the windows overlap, their counts differ by length, and the
length with the highest score was identified on this cohort.
\end{abstract}

\section{Study design}
\subsection{Cohort}
The source is a frozen 166-session Sydney snapshot. Bus, car, and train
recordings with snapshot duration of at least 240 seconds are candidates;
there is no upper-duration limit. Event-data checks and feature-quality rules
require each included session to have at least one valid window at every
tested length. The resulting common cohort has 149 sessions from seven
participants: 38 bus, 66 car, and 45 train. Session identities and
participant folds remain the same across lengths. Selection for long,
quality-surviving sessions limits the population to which these estimates
apply.

\subsection{Windows and features}
Window lengths are 20, 40, \ldots, 240 seconds. Each advances by half its
length, giving 50\% overlap within a session. For a window of length $W$
seconds, the feature pipeline requires at least $30W$ accelerometer, $4W$
gyroscope, and $2W$ magnetometer samples. A required sensor stream with a
consecutive-sample gap over 500 ms causes rejection. The fixed 11-feature
vector contains mean, standard deviation, range, and minimum of
accelerometer and gyroscope magnitudes, plus standard deviation, range, and
minimum of magnetometer magnitude. Sensor inputs, feature definitions,
quality rules, and overlap fraction are held fixed across lengths.

\subsection{Training and scoring}
For each length, leave-one-participant-out evaluation trains on six
participants and tests on the seventh. The XGBoost hyperparameters were
selected on a separate NOR-TMD dataset and remain fixed here. Each fold is
trained only on Sydney windows, using class-balanced training weights.
Seeds 42, 43, and 44 give 36 length-and-seed evaluations and 252 model fits
across the seven folds. No held-out participant contributes training windows
to its fold.

Within each length and seed, predictions from all seven held-out folds are
pooled. Macro F1 is the mean of bus, car, and train F1; balanced accuracy is
the mean of their recalls. Reported scores are means of the three seed-level
pooled metrics. Each session is tested once per seed and length, but shorter
lengths produce more overlapping windows. Pooled scores therefore weight
sessions by their retained window counts. They are neither participant-equal
scores nor measures from independent journeys.

\section{Results}
Mean macro F1 rises from 0.596 at 20 seconds to 0.693 at 120 seconds, then
varies from 0.662 to 0.682 across the remaining lengths. Balanced accuracy
follows the same broad pattern and reaches 0.691 at 120 seconds. Table~\ref{tab:window-grid}
shows the complete grid, including the decline in retained window counts as
length increases.

\begin{table}[htbp]
\centering
\small
\caption{Mean pooled held-out scores across three seeds on the common
149-session cohort. Window counts are also means across seeds.}
\label{tab:window-grid}
\begin{tabular}{rrrr}
\hline
Window (s) & Mean windows & Macro F1 & Balanced accuracy \\
\hline
20  & 26,422 & 0.596 & 0.599 \\
40  & 13,099 & 0.625 & 0.626 \\
60  & 8,657  & 0.643 & 0.643 \\
80  & 6,435  & 0.652 & 0.652 \\
100 & 5,107  & 0.668 & 0.666 \\
120 & 4,217  & 0.693 & 0.691 \\
140 & 3,580  & 0.682 & 0.679 \\
160 & 3,105  & 0.669 & 0.667 \\
180 & 2,733  & 0.675 & 0.673 \\
200 & 2,439  & 0.669 & 0.667 \\
220 & 2,196  & 0.665 & 0.663 \\
240 & 1,997  & 0.662 & 0.660 \\
\hline
\end{tabular}
\end{table}

Table~\ref{tab:window-classes} shows class F1 at the shortest length, the
observed peak, and the longest length. Bus is the most difficult class at
120 seconds despite its improvement across the grid. Train F1 decreases at
the longest length. These patterns are descriptive and do not establish
which aspect of longer temporal context caused the changes.

\begin{table}[htbp]
\centering
\small
\caption{Mean class F1 across three seeds at selected window lengths.}
\label{tab:window-classes}
\begin{tabular}{rrrr}
\hline
Window (s) & Bus F1 & Car F1 & Train F1 \\
\hline
20  & 0.358 & 0.745 & 0.684 \\
120 & 0.540 & 0.818 & 0.722 \\
240 & 0.513 & 0.808 & 0.664 \\
\hline
\end{tabular}
\end{table}

\section{Interpretation and limitations}
The 120-second length has the highest observed mean macro F1 in this grid.
It requires 120 seconds of sensor history for a first full-window prediction
and advances in 60-second steps. No on-device latency, energy use, or
transition-detection delay was measured, so the operational trade-off cannot
be quantified from these scores.

The seven participants, class concentration, shared sessions across lengths,
and window-level pooling limit generalization beyond this cohort. Feature
definitions were informed by earlier Sydney exploration, although model
hyperparameters were selected on separate data. The peak was selected after
examining all 12 lengths; it should not be treated as a confirmed optimum.
An independent participant cohort with the length fixed in advance is needed
to assess whether the observed peak persists.

\end{document}
